\documentclass[11pt]{article}
\usepackage[margin=1in]{geometry}
\usepackage[T1]{fontenc}
\usepackage{amsmath,amssymb,mathtools,bbm}
\usepackage{booktabs,longtable,array,enumitem}
\usepackage{listings,xcolor,tikz,hyperref}
\sloppy
\hypersetup{colorlinks=true, linkcolor=blue!50!black, urlcolor=blue!50!black}
\definecolor{cbg}{rgb}{0.97,0.97,0.97}
\lstset{language=Python, basicstyle=\ttfamily\footnotesize,
  backgroundcolor=\color{cbg}, showstringspaces=false, breaklines=true,
  frame=single, framerule=0pt, numbers=left, numberstyle=\tiny, xleftmargin=2em}
\newcommand{\code}[1]{\texttt{#1}}
\newcommand{\file}[1]{\texttt{#1}}
\newcommand{\sbscode}[1]{\begin{minipage}[t]{0.55\textwidth}#1\end{minipage}}
\newcommand{\sbstext}[1]{\hfill\begin{minipage}[t]{0.42\textwidth}\small #1\end{minipage}}

\title{AutoBayes milestone M0: evidence for review}
\author{Claude session, branch \code{m0/gaussian-varying-intercept}}
\date{September 8, 2026}

\begin{document}
\maketitle

\section{What M0 delivers}

One executable path for the Gaussian varying-intercept model
\[
y_i \sim \mathrm{Normal}(\mu_i,\sigma), \qquad
\mu_i = a_c + (x_i - \bar x)\,b + u[g_i], \qquad u[j] = s_g\, z_j,\quad z_j \sim \mathrm{Normal}(0,1),
\]
with explicit priors on the four brms classes (\code{Intercept}, \code{b}, \code{sd}, \code{sigma}),
an inspectable PyMC model, posterior and predictive draws, PyMC's NUTS on the CPU, and a
JAX/NumPyro path with an explicit device choice. Everything outside that family is
refused with an error naming the construct.

\begin{center}\footnotesize
\begin{tabular}{p{0.2\textwidth}p{0.74\textwidth}}
\toprule
Branch, clone & \code{m0/gaussian-varying-intercept} in \file{IQSS-Research/projects/AutoBayes}, five commits on \code{f87f639} \\
Package & \file{autobayes/} (\file{formula.py}, \file{priors.py}, \file{model.py}, \file{inference.py}, \file{predict.py}, \file{diagnostics.py}) \\
Tests & \file{tests/}: 36 fast, 2 slow (\code{-m slow}); all 38 pass \\
Reference density & \file{tests/reference\_logp.py}, imports nothing from \code{autobayes}, PyMC, or Bambi \\
Example & \file{examples/m0\_varying\_intercept.py}; outputs captured in \file{docs/m0-evidence/} \\
Decision record & \file{docs/m0-decision.md} (written before implementation), \file{docs/coordination/STATUS.md} \\
Environment & \file{.venv}: pymc 5.28.5, bambi 0.17.2, numpyro 0.21.0, jax 0.10.2 (CPU wheel), arviz 0.23.4, pytensor 2.38.3 \\
\bottomrule
\end{tabular}
\end{center}

\section{Bambi already builds this model; the layer is thin on purpose}

The milestone asked for a comparison before any new code. Verified against the installed
Bambi 0.17.2: it parses \code{y \textasciitilde{} x + (1|g)}, builds the one-hot group matrix with levels
in sorted order, accepts explicit priors including a hyperprior on the group SD, uses the
non-centred parameterisation by default (as brms always does), centres the predictors by
default and back-transforms the intercept after sampling (as brms does), exposes the PyMC
model as \code{model.backend.model} after \code{build()}, forwards \code{**kwargs} to
\code{pm.sample} so \code{nuts\_sampler="numpyro"} with \code{nuts\_sampler\_kwargs=\{"chain\_method": ...\}}
reaches PyMC 5.28.5's \code{sample\_jax\_nuts}, and predicts for new group levels with brms's
default \code{sample\_new\_levels="uncertainty"} semantics only.

The layer therefore adds six things, each because Bambi does not do it: (1)~a grammar
check that refuses \code{(1 + x | g)}, \code{(1 | g/h)}, two grouping factors,
interactions, transformations, \code{0 +}, categorical predictors, and missing data with a
message naming the construct; (2)~mandatory explicit priors by class, with no silent
defaults; (3)~a \code{describe()} that states the density actually sampled, centring and
transforms included; (4)~device resolution that raises rather than falling back, enables
float64, sets the host device count for parallel CPU chains, and reports the actual
device; (5)~predictive draws with named new-group semantics, including the population
option (brms \code{"gaussian"}) that Bambi lacks; (6)~the diagnostic panel with the
thresholds from \file{reference/diagnostics.md}. No parser, design-matrix builder, prior
language, or sampler was written.

\section{The model that is actually sampled}

\code{model.describe()} on the example data prints the specification below; the PyMC
representation follows it. Both are real output (\file{example\_pymc\_cpu.txt},
\file{pymc\_model\_repr.txt}).

\begin{lstlisting}[numbers=none]
Linear predictor
  mu_i = a_c + (x_x - mean(x_x)) * b_x + u[school_i]
  Predictors are centred at their sample means (center_predictors=True, as in
  brms and Bambi). a_c is the expected response at the predictor means; the
  prior on the intercept applies to a_c. The posterior reports the raw-scale
  intercept a = a_c - sum_p mean(x_p) * b_p, computed after sampling.
Group effects (non-centred)
  u[j] = sd_school * z[j],  z[j] ~ Normal(0, 1),  j = 0..19
PyMC variables sampled (value space)
  sigma_log__              log(sigma)
  Intercept                a_c
  x                        b_x
  1|school_sigma_log__     log(sd_school)
  1|school_offset          z
\end{lstlisting}
\begin{lstlisting}[numbers=none]
     sigma ~ HalfNormal(0, 1)          1|g_offset ~ Normal(0, 1)
 Intercept ~ Normal(0, 5)                     1|g ~ Deterministic(f(1|g_offset, 1|g_sigma))
         x ~ Normal(0, 2)                      mu ~ Deterministic(f(Intercept, 1|g_offset, x, 1|g_sigma))
 1|g_sigma ~ HalfNormal(0, 1)                   y ~ Normal(mu, sigma)
\end{lstlisting}

Two facts a researcher must know and would not see in Bambi's own printout. First, the
\code{Intercept} in the PyMC model is $a_c$, the intercept at the predictor means; the
\code{Intercept} in the returned trace is $a = a_c - \bar x\, b$. The prior
$\mathrm{Normal}(0,5)$ is on $a_c$. With the example data $\bar x = -0.092$ and $b\approx 0.49$
the two differ by $0.045$, small here and not small in general. Second, the sampler moves
$\log\sigma$ and $\log s_g$, not $\sigma$ and $s_g$; the joint density in that space carries
the Jacobian $+\log\sigma + \log s_g$, which is what the reference below reproduces.

\section{Log density and meaning (checks A1--A7)}

\begin{minipage}[t]{0.54\textwidth}\begin{lstlisting}[firstnumber=65]
    y, Xc, _, gidx, levels = build_design(data, response, predictors, group, center)
    a, b = float(point["a_c"]), np.atleast_1d(point["b"])
    log_sigma, log_sd = float(point["log_sigma"]), float(point["log_sd"])
    z = np.asarray(point["z"], dtype=float)
    assert z.shape == (len(levels),)
    sigma, sd = np.exp(log_sigma), np.exp(log_sd)

    mu = a + Xc @ b + sd * z[gidx]
    lp = np.sum(stats.norm.logpdf(y, loc=mu, scale=sigma))
    lp += np.sum(stats.norm.logpdf(z, 0.0, 1.0))
    lp += prior_logpdf(*priors["intercept"], a)
    for i, p in enumerate(predictors):
        lp += prior_logpdf(*priors["b"][p], b[i])
    lp += prior_logpdf(*priors["sd"], sd)
    lp += prior_logpdf(*priors["sigma"], sigma)
    if jacobian:
        lp += log_sigma + log_sd
\end{lstlisting}\end{minipage}
\sbstext{\file{tests/reference\_logp.py:65--81}. The reference builds its own centred
design and its own sorted level index from the data frame, evaluates every prior with
SciPy (half-Student-t written out), and adds the log-transform Jacobian. It shares no code
with the library, so agreement is evidence about Bambi's graph, not about the layer.

The comparison point is a dictionary in PyMC's value space produced by
\code{Model.value\_point()}, whose keys are checked against \code{pymc\_model.value\_vars}.}

\medskip
\begin{center}\footnotesize
\begin{tabular}{lp{0.6\textwidth}ll}
\toprule
Check & What is compared & Tolerance & Result \\
\midrule
A1 & \code{compile\_logp()} vs reference, 5 random points, 4 formula/centring variants & $10^{-8}$ abs & pass \\
A2 & same with \code{jacobian=False}; difference equals $\log\sigma+\log s_g$ & $10^{-8}$, $10^{-12}$ & pass \\
A3 & \code{compile\_dlogp()} vs central differences of the reference, 3 points & rel $10^{-5}$ & pass \\
A4 & JAX-compiled density (\code{get\_jaxified\_logp}) vs PyTensor path & $10^{-8}$ abs & pass \\
A5 & scales fixed by \code{pm.do}: negative Hessian equals closed-form precision at two points; maximiser equals exact posterior mean & $10^{-8}$, $10^{-4}$ & pass \\
A6 & $X_\text{model} = x-\bar x$, $Z$ one-hot in sorted level order, \code{mu} and \code{1|g} from the graph vs by hand & $10^{-12}$ & pass \\
A7 & 17 unsupported formulas and 5 data/prior faults raise the named error & --- & pass \\
\bottomrule
\end{tabular}
\end{center}

Three PyMC 5.28.5 facts surfaced while making A4--A6 pass, all verified in source and
none a model problem. \code{get\_jaxified\_logp(model, negative\_logp=...)} has its flag
named backwards: the default \code{True} returns the log density and \code{False} negates
it. \code{compile\_d2logp()} returns the \emph{negative} Hessian
(\code{pymc.pytensorf.hessian(negate\_output=True)}), i.e. the precision matrix, which is
why A5 compares it to $+P$. A deterministic such as \code{mu} is a graph of random
variables; to evaluate it at a point in value space the RVs must first be substituted
(\code{model.replace\_rvs\_by\_values}). Each is recorded in the test that hit it.

\section{Sampling (checks B1--B2)}

\paragraph{B1: exact conditional posterior.} With $\sigma = 0.3$ and $s_g = 0.5$ fixed
through \code{pm.do} on the inspectable model, the posterior of $(a_c, b, z_1..z_6)$ is
Gaussian with precision $P = \mathrm{diag}(1/25, 1/4, 1,\dots,1) + W^\top W/\sigma^2$,
$W = [\mathbf 1,\ x-\bar x,\ s_g Z]$. NUTS on that model (n = 60, J = 6) against the exact
mean, with MCSE $= \text{exact SD}/\sqrt{\mathrm{ESS}}$:

\begin{center}\footnotesize
\begin{tabular}{lrrrrrr}
\toprule
\multicolumn{7}{l}{\emph{Committed test: seed 11, 4 chains $\times$ 2000 draws, threshold $|z|<4$ and SD within 10\%}} \\
 & exact mean & MCMC mean & exact SD & MCMC SD & ESS & $z$ \\
\midrule
$a_c$ & 1.2147 & 1.2289 & 0.2078 & 0.2149 & 1273 & $+2.45$ \\
$b_x$ & 0.4561 & 0.4565 & 0.0423 & 0.0423 & 3159 & $+0.48$ \\
$z_0 \dots z_5$ & & & 0.437--0.456 & 0.450--0.467 & 1413--1503 & $-2.20$ to $-2.80$ \\
\bottomrule
\end{tabular}

\medskip
\begin{tabular}{llrrrrrr}
\toprule
\multicolumn{8}{l}{\emph{Seed check, 4 chains, 2000 tune; $z$ for $a_c$, for the identified direction $a_c + s_g\bar z$, and for the weak direction $a_c - s_g\bar z$}} \\
seed & draws & ESS $a_c$ & ESS ident. & ESS weak & $z(a_c)$ & $z$(ident.) & $z$(weak) \\
\midrule
11 & 2000 & 1315 & 8929 & 1294 & $+0.06$ & $-0.96$ & $+0.09$ \\
12 & 4000 & 2493 & 17328 & 2444 & $-1.69$ & $+1.72$ & $-1.76$ \\
13 & 4000 & 2474 & 17622 & 2429 & $+1.59$ & $+0.89$ & $+1.57$ \\
14 & 8000 & 4884 & 31828 & 4793 & $-1.14$ & $-0.15$ & $-1.13$ \\
\bottomrule
\end{tabular}
\end{center}

The committed run passed its threshold ($|z| < 4$, SDs within 10\%) but its eight
$z$-scores share one sign pattern, so they are one event, not eight. The pattern is the
direction $a_c - s_g \bar z$: the data pin down $a_c + s_g \bar z$ (the grand mean, exact
SD 0.040, about one effective draw per draw) and leave $a_c - s_g \bar z$ to the priors
(exact SD 0.41, about one effective draw per six). Across seeds and lengths the deviations
alternate sign and stay under 1.8 MCSE, so the first run was a $2.5\sigma$ draw and not a
bias. Files: \file{B1\_conditional\_gaussian.json}, \file{B1\_seed\_check.json}.

\paragraph{B2: simulated recovery on both backends.} $n = 200$, $J = 20$, truth
$a = 1.0$, $b = 0.5$, $s_g = 0.5$, $\sigma = 0.3$; 4 chains, 1000 tune, 2500 draws,
\code{target\_accept} 0.9, seed 7.

\begin{center}\footnotesize
\begin{tabular}{lrrrrrr}
\toprule
 & mean & SD & 94\% HDI & $\hat R$ & bulk ESS & tail ESS \\
\midrule
\multicolumn{7}{l}{\emph{PyMC NUTS, CPU: 0 divergences of 10000, E-BFMI 0.70, tree-depth saturation 0\%}} \\
Intercept & 0.776 & 0.143 & [0.491, 1.037] & 1.003 & 755 & 1384 \\
$b_x$ & 0.486 & 0.024 & [0.441, 0.529] & 1.000 & 5055 & 5721 \\
$s_g$ & 0.642 & 0.113 & [0.450, 0.858] & 1.003 & 1215 & 2544 \\
$\sigma$ & 0.306 & 0.016 & [0.276, 0.336] & 1.001 & 5208 & 6235 \\
\midrule
\multicolumn{7}{l}{\emph{NumPyro NUTS via JAX, CPU: 0 divergences of 10000, E-BFMI 0.67, tree-depth saturation 0\%}} \\
Intercept & 0.779 & 0.146 & [0.518, 1.082] & 1.005 & 771 & 1235 \\
$b_x$ & 0.487 & 0.024 & [0.442, 0.531] & 1.001 & 5401 & 6061 \\
$s_g$ & 0.647 & 0.117 & [0.443, 0.862] & 1.008 & 1010 & 1966 \\
$\sigma$ & 0.305 & 0.017 & [0.274, 0.337] & 1.000 & 5229 & 6127 \\
\bottomrule
\end{tabular}
\end{center}

Every panel row passes on both backends and the truth is inside every 94\% HDI. The two
backends' posterior means differ by 0.0027 (Intercept), 0.0009 ($b$), 0.0052 ($s_g$),
0.0004 ($\sigma$), against 4-combined-MCSE tolerances of 0.030, 0.0019, 0.019, 0.0013:
the same posterior, sampled two ways.

Two readings that would be wrong. The Intercept posterior sits at 0.78 against a nominal
1.0 and $s_g$ at 0.64 against 0.5: this is not shrinkage error. The twenty simulated
effects have realised mean $-0.234$ and realised SD 0.602, so the quantities the data can
identify are $1.0 - 0.234 = 0.766$ and about 0.60, both near the posterior centres. And at
4 chains $\times$ 1000 draws these same fits gave $\hat R = 1.02$ and bulk ESS 350 on the
Intercept on both backends, for the reason B1 exhibits: the weakly identified
intercept-versus-mean-offset direction. The runs were lengthened; the thresholds were not
loosened. One recovery on one simulated data set is not a calibration claim.

\section{Devices: executed versus supported}

\begin{center}\footnotesize
\begin{tabular}{p{0.29\textwidth}p{0.27\textwidth}p{0.35\textwidth}}
\toprule
Path & Status & Evidence \\
\midrule
\code{backend="pymc"}, CPU & executed & B2, \file{example\_pymc\_cpu.txt} (wall 5.0 s incl.\ compile) \\
\code{backend="numpyro"}, \code{device="cpu"} & executed, 4 parallel CPU chains, float64 & B2, \file{example\_numpyro\_cpu.txt} (3.9 s) \\
\code{backend="numpyro"}, \code{device="gpu"} & supported in code, \textbf{not validated} & \file{example\_numpyro\_gpu\_attempt.txt}: the refusal \\
\code{backend="pymc"}, \code{device="gpu"} & refused by design & \code{test\_gpu\_on\_pymc\_backend\_is\_error} \\
\bottomrule
\end{tabular}
\end{center}

\begin{minipage}[t]{0.54\textwidth}\begin{lstlisting}[firstnumber=70]
    if device == "cpu":
        # Must run before the JAX backend initialises, else it is a no-op.
        numpyro.set_host_device_count(int(chains))
    jax.config.update("jax_enable_x64", True)
    devs = jax.devices()
    if device == "gpu":
        gpus = [d for d in devs if d.platform in _GPU_PLATFORMS]
        if not gpus:
            raise DeviceError(
                "device='gpu' requested for backend='numpyro', but JAX sees no GPU. "
                f"Visible devices: {_jax_devices_summary(jax)}. Install a CUDA-enabled "
                "jaxlib on a machine with an NVIDIA GPU, or pass device='cpu'. "
                "AutoBayes does not fall back to CPU silently."
            )
        return gpus[0], "vectorized"
\end{lstlisting}\end{minipage}
\sbstext{\file{autobayes/inference.py:70--84}. The report on every fit records the device
the arrays actually landed on (probed inside \code{jax.default\_device}), the dtype, and
the chain method; on this machine it reads \code{cpu:cpu(id=0)}, \code{float64},
\code{parallel}. The two wall-clock figures above are not a speed comparison: same
machine, same targets, but compilation is inside both numbers and the sizes are tiny.}

\medskip
The GPU command is in \file{docs/m0-gpu-validation.md} and \file{cluster/m0\_gpu\_smoke.sbatch};
the cluster's existing \code{electit\_env2} lacks Bambi, so it needs a new environment.
Nothing was submitted, per the handoff. What a GPU run would add is one fact: that the same
compiled graph runs on the other device. It would not add a speed claim.

\section{Prediction meaning (check A8)}

In-sample \code{mu} from \code{predict()} equals $a + x b + u[g]$ computed by hand from
the trace to $10^{-10}$; \code{y} draws scatter around it with pooled SD equal to the
posterior mean of $\sigma$. For an unseen group, \code{new\_groups="population"} draws
$u_\text{new} \sim \mathrm{Normal}(0, s_g)$ per posterior draw, so its predictive SD is
about $\sqrt{\sigma^2 + s_g^2}$: in the example 0.74 for a new school against 0.32 for an
observed one ($\sqrt{0.306^2 + 0.642^2} = 0.71$ before posterior spread in $a$ and $s_g$).
\code{"resample\_observed"} reproduces Bambi's behaviour, verified draw by draw to be one of
the observed levels' effects; it is brms's default and it understates between-group
variance, which is why it is not the default here.

\section{Not claimed, and next}

Not claimed: any family beyond Gaussian, varying slopes, multiple or nested grouping
factors, categorical predictors, missing data, calibration in general, GPU execution, or
any speedup. Follow-ups worth a decision: run \file{cluster/m0\_gpu\_smoke.sbatch} once a
Bambi-capable environment exists on the cluster; decide whether the three PyMC quirks
belong in \file{reference/libraries.md} (the pack is a fixed input, so that is Andrew's
call); and choose the next family (M1) only after this review.

\end{document}
